\section{Experimental Framework}\label{sec:framework}

This section defines the evaluation platform, metrics, and model configurations that underpin the analyses in Sections~\ref{sec:fake_safety}--\ref{sec:avoidance_source}. Our primary methodological contribution is the displacement-filtered Collision Avoidance Rate (dfCAR), which separates genuine avoidance from action collapse.

\subsection{Benchmark and Evaluation Protocol}

We conduct all experiments on SafeLIBERO~\cite{hu2025vlsa}, a collision-aware evaluation extension of the LIBERO simulation platform~\cite{liu2023libero}. SafeLIBERO augments standard manipulation benchmarks with per-timestep collision detection and non-target obstacle annotations, enabling systematic measurement of both task completion and safety.

We select the \textit{object\_pudding} Level~I suite, which comprises four tabletop manipulation tasks. In each task, a target object is surrounded by non-target obstacles that the robot must avoid during grasping and transport. The MuJoCo physics engine provides collision detection: any contact between the robot (end-effector or arm links) and a non-target object registers as a collision event.

Each method is evaluated for 50 episodes per task (200 episodes total). For multi-seed experiments (Pi0.5), we run 3 seeds to produce 600 episodes per configuration. Initial object configurations are fixed across methods to ensure fair comparison.

\subsection{Evaluation Metrics}\label{sec:metrics}

\paragraph{Standard metrics.} We report two conventional metrics. The \emph{Task Success Rate} (TSR) is the fraction of episodes where the manipulation task reaches completion. The \emph{Collision Avoidance Rate} (CAR) is the fraction of episodes in which no collision occurs between the robot and any non-target object.

\paragraph{Displacement-filtered CAR.} Standard CAR conflates two qualitatively distinct behaviors: a robot that plans a trajectory around obstacles, and a robot that has collapsed into inaction. Both register as ``safe,'' yet only the former reflects collision avoidance capability. We propose the displacement-filtered CAR (dfCAR), which excludes episodes exhibiting action collapse from both numerator and denominator.

We define an episode $e$ as \emph{action-collapsed} when its maximum end-effector displacement satisfies $d_{\max}(e) < \tau_d$, where $\tau_d = 0.1$\,mm. The filtered metric is then:
\begin{equation}\label{eq:car_f}
\text{dfCAR} = \frac{|\{e : \text{safe}(e) \;\land\; d_{\max}(e) \geq \tau_d\}|}{|\{e : d_{\max}(e) \geq \tau_d\}|}
\end{equation}
where $d_{\max}(e)$ denotes the maximum Euclidean displacement of the end-effector across all timesteps in episode $e$. The threshold $\tau_d = 0.1$\,mm corresponds to sub-actuator-resolution motion. In practice, the displacement distribution is bimodal: collapsed episodes cluster below 1\,$\mu$m and active episodes exceed 200\,mm, with no intermediate population. Consequently, any $\tau_d \in [2, 200]$\,mm produces identical filtering (Appendix~\ref{sec:appendix_threshold}).

\paragraph{Auxiliary metrics.} We additionally report the \emph{Action-Collapse Rate}, the fraction of episodes with $d_{\max} < \tau_d$, which quantifies the severity of action collapse for each method. The \emph{Simultaneous Success Rate} (SSR $= $ TSR $\cap$ CAR) measures the fraction of episodes that achieve both task completion and collision avoidance.

\subsection{Models and Training}

\paragraph{Base architectures.} We evaluate two VLA models from the Pi0 family~\cite{black2024pi0}. \textbf{Pi0} (${\sim}$3B parameters) uses a PaliGemma vision-language encoder paired with a dedicated action expert that generates continuous actions via flow matching~\cite{lipman2023flow}. \textbf{Pi0.5} extends this architecture to ${\sim}$3.5B parameters with broader pre-training data, and serves as a stronger baseline that exhibits less action collapse.

\paragraph{Training configurations.} We evaluate four model variants:
\begin{itemize}
\item \textit{Full-FT baseline}: Pi0 with all parameters fine-tuned on SafeLIBERO demonstrations using the flow-matching action loss.
\item \textit{LoRA baseline}: Pi0 with LoRA adapters (${\sim}$33.5M trainable parameters), same demonstrations and loss.
\item \textit{A1-LoRA}: Pi0 + LoRA augmented with two auxiliary heads. An SDF Decoder (4-layer MLP, ${\sim}$395K parameters) predicts the signed distance from the end-effector to the nearest obstacle surface, supervised with a clamped $L_1$ loss and an Eikonal regularizer ($\mathcal{L}_\text{SDF}$). An EE Extractor (2-layer MLP, ${\sim}$263K parameters) predicts end-effector position from action expert features ($\mathcal{L}_\text{EE}$, MSE loss). SDF ground truth is computed from known obstacle meshes in simulation.
\item \textit{Pi0.5 + A1}: The A1 architecture applied to the Pi0.5 backbone. Training proceeds in two phases: Phase~1 (steps 0--5K) trains only the auxiliary heads on detached features with the backbone frozen; Phase~2 (steps 5K--30K) unfreezes LoRA parameters for joint optimization. We evaluate $\lambda_\text{SDF} \in \{0.005, 0.01\}$.
\end{itemize}

\subsection{Oracle SDF Refinement}

To test the theoretical upper bound of end-effector-only SDF refinement (analyzed in Section~\ref{sec:ee_only}), we replace the learned SDF network with the ground-truth signed distance field computed from simulation geometry. This eliminates prediction error and isolates the structural limitation of the refinement paradigm itself.

We evaluate three oracle variants that differ in how the SDF gradient is applied:
\begin{itemize}
\item \textit{Oracle-adaptive}: runtime gradient refinement with a dynamic step size $\alpha$ that scales inversely with gradient magnitude.
\item \textit{Oracle-fixed}: a constant step size $\alpha = 0.005$.
\item \textit{Oracle-feature}: gradient applied in the VLA's latent feature space rather than Cartesian action space.
\end{itemize}
Each variant is evaluated for 160 episodes (4 tasks $\times$ 40 episodes) with a single seed.

\subsection{SDF Ground Truth}

The workspace signed distance field is computed from the known obstacle meshes available in the LIBERO/MuJoCo simulation environment. SDF values are positive in free space, zero on the obstacle surface, and negative inside the obstacle volume. Ground-truth labels are generated at query points sampled along predicted end-effector trajectories during training.